\scanpage{14}
\mspage{12}
les (\unsure{petites ?|petites}) catégories, et \deleted{les \unsure{fl.}} où on définit
\[
\operatorname{Hom}_{\mathrm{Cathot}}(A,B)\overset{\mathrm{déf}}{=}
\underbrace{\pi_0(\underline{\mathrm{Hom}}(A,B))}_{\substack{\text{classes d'homotopie}\\\text{de foncteurs de }A\\\text{dans }B}}
\]
C'est cette conclusion qui sera pour nous utile, et non la description explicite du chemin $h$. Cette existence admise (cf ci-dessous la preuve), on voit donc que si on regarde, \inserted{pour $a,b\in\operatorname{Ob}X$, \deleted{\unsure{$\nu\in$}}} le système des catégories $\deleted{\unsure{\mathrm{Hom}}}\underline{\mathrm{Ch}}_{\Sigma,\unsure{\nu|n}}(X;a,b)$ \inserted{$(\nu\in\mathbb N)$} comme un \uline{ind-objet de $\mathrm{Cathot}$}, qu'on \deleted{\unsure{va}} désignera par
\[
\begin{aligned}
&\deleted{\unsure{\Phi_n C_\Sigma}}\\
&\underline{\mathrm{Ch}}_\Sigma^{\mathrm{ind}}(X;a,b)\in\operatorname{Ob}\operatorname{Ind}(\mathrm{Cathot}),
\end{aligned}
\]
on aura un \deleted{application \unsure{dans la}} \inserted{\deleted{\unsure{f}} morphisme} canonique dans $\operatorname{Ind}(\mathrm{Cathot})$
\[
\underline{\mathrm{Ch}}_\Sigma^{\mathrm{ind}}(X;x,y)\times\underline{\mathrm{Ch}}_\Sigma^{\mathrm{ind}}(X;y,z)\xrightarrow{*}\underline{\mathrm{Ch}}_\Sigma^{\mathrm{ind}}(X;x,z),
\]
appelé ``\uline{\deleted{\unsure{morphisme}} \inserted{ind-morphisme} de composition des chemins}''.

On aurait aimé le déduire d'un honnête \uline{foncteur} ``composition des chemins''
\[
\deleted{\unsure{\Phi}}\underline{\mathrm{Ch}}_\Sigma(X;x,y)\mathbin{\unsure{*|\times}}\underline{\mathrm{Ch}}_\Sigma(X;y,z)\xrightarrow{?}\underline{\mathrm{Ch}}_\Sigma(X;x,z),
\]
qu'on noterait $(c,c')\longmapsto c'\circ c$, ou $\unsure{c'c|c',c}$, mais je doute qu'un tel foncteur existe,

% archive sheet number: 13
